\section{Threats to Validity and Limitations}
\label{sec:threats}

We hold this work to the standard of self-criticism it advocates, and enumerate the threats in
descending order of how much they could change the paper's claims. We distinguish \emph{threats to
the thesis} (is the gain of training-free RL really governed by reward spread, and does that buy us
anything at the agent level?) from \emph{threats to the moat} (is the contribution durable and
novel?) and \emph{threats to the empirics} (do the measurements support the conclusions?). The most
consequential item is not an external risk but an internal, honest negative result of our own, so we
state it first.

\paragraph{Our strongest finding is a negative one: paralinguistic agentic recovery is, on our own
data, unsupported.}
The constructive reading of OSA-2 depends entirely on \Cref{conj:spreadfloor}: that the frozen-model
context-isolation construction actually yields \emph{non-degenerate} per-block verifiable rewards,
i.e.\ a per-component spread $\spread_i$ bounded away from zero so that each $\gain(\Rwd_i)>0$. This
is an empirical conjecture, not a theorem, and our own measurements indicate it \emph{fails or is
fragile precisely for the flagship's paralinguistic factors}. For speaker, the kNN probe sits at
$\approx 0.04$ --- near chance --- at \emph{all} $37$ hidden layers and under both pooling readouts,
so the frozen omni embedding affords essentially no speaker spread to optimize over. For emotion, a
rigorous multi-seed re-run replaces the earlier single-seed headline: with the readout layer selected
on dev, a paired-delta bootstrap, and five seeds ($42, 7, 123, 2024, 31337$), the mean improvement is
$\Delta = +0.037$ (range $[-0.053, +0.097]$), and the paired CI excludes zero in only $2$ of $5$
seeds. The previously reported single-seed $+0.097$ was the \emph{best} seed under \emph{oracle
test-layer} selection --- a winner's-curse / selection-leak artifact, not a stable effect. The
committed artifact is
\texttt{projects/speech-mllm-omni-embedding-rl/\_repro/emotion\_pool\_paired\_v2.json}. The honest
conclusion is that the agentic-recovery move is empirically supported for \emph{content/intent}
factors (where single-model task-level selection already realizes large, well-separated spread ---
e.g.\ SLURP $+0.330$, \Cref{sec:feasibility}) but \emph{not} for paralinguistics, where on our data
\Cref{conj:spreadfloor} is likely false. We present this as a central finding of the paper, not as a
failure to be hidden: it sharply scopes where the reward-spread thesis has empirical teeth.

\paragraph{Context isolation buys nothing structurally: the isolated optimum equals the monolithic one.}
A reader might suppose that decomposing a task into context-isolated blocks is itself the source of
gain. \Cref{thm:qstarprod} (\texttt{qstar\_product}) forbids that reading: under a product base and a
separable reward $\Rwd=\sum_i\Rwd_i$, the global Gibbs optimum factorizes as
$\qs=\prod_i q^*_i$, so tilting each isolated block to its own optimum assembles \emph{exactly} the
same distribution --- and the same objective value $\Fobj{\qs}=\beta\log\Zpart$ --- as a single
monolithic tilt against $\Rwd$. Isolation \emph{per se} adds zero gain. The only way to increase
$\gain$ is to introduce additional \emph{non-degenerate} verifiable rewards $\Rwd_i$ (each strictly
positive by \texttt{gain\_pos\_of\_nonconstant}); the architectural act of separating contexts is a
convenience for credit assignment, not a source of optimization headroom. Any claim that the agent
structure realizes gain a monolithic search could not is therefore an \emph{empirical} claim about
the availability of new non-degenerate rewards --- exactly the claim the previous paragraph shows is
unsupported for paralinguistics.

\paragraph{The moat is an absence of evidence (fast-follower risk).}
Our central feasibility argument rests on the observation that no prior work couples training-free,
reward-guided inference-time optimization with a frozen omni model's disentangled embeddings as the
search variable. An absence of evidence is structurally weaker than evidence of absence: it is
falsified the moment a competing group publishes the same coupling, and the individual ingredients
--- best-of-$N$ and reranking \citep{beirami2024bon,verdun2025softbon}, controlled decoding
\citep{mudgal2024controlled}, no-gradient credit assignment \citep{jitrl2026}, test-time RL
\citep{zuo2025ttrl,li2025tpo} --- are all public. We therefore do not claim a defensible technical
monopoly; we claim a \emph{first-mover synthesis} whose value is the formalization (the OSA theorems)
and the machine-checked qualitative core, not the mechanism itself. A reader should weight this
contribution as integrative rather than as a primitive discovery.

\paragraph{The mechanism is not novel; it is a domain transfer with a now-doubtful precondition.}
The Gibbs-optimal tilting $\qs \propto \qo\exp(\Rwd/\beta)$ and its KL-regularized variational
characterization are textbook results in the alignment-as-Bayesian-inference literature
\citep{korbak2022kl,yang2024asymptotics}, and the best-of-$N$/KL frontier is now sharply understood
\citep{beirami2024bon,gui2024bonbon}. Our contribution is to instantiate this machinery over the
\emph{embedding} action space $\cZA$ of a frozen speech model rather than the token space. The
transfer assumes the relevant degrees of freedom $\cZ$ already carry the pretrained structure the
reward asks for \citep{frozenllmparaling2024,salmonn2023}. Our results split this assumption: the
\emph{semantic/content} structure is clearly present and high-spread, but the \emph{paralinguistic}
structure (speaker, emotion) is, on our probes, at or near chance --- so for those factors the hidden
precondition of the transfer is violated, which is exactly why the negative result above is not
incidental but predicted by the theory.

\paragraph{OSA-2 idealizes context isolation as exact separability.}
OSA-2 models a multi-key agent reward as a sum of per-context-isolated components and proves
additivity under \emph{exact} separability of $\Rwd$ across the content, speaker, emotion, and
language/intent factors. Real agent rewards are at best \emph{approximately} separable: an
ASR-conditioned downstream reward leaks into the emotion channel through prosody, and intent
classification is not independent of transcription quality. Under a residual coupling of magnitude
$\epsilon$ in the cross-factor Hessian, the Gibbs optimum of the separable surrogate is displaced from
the true optimum, and the additive $\gain$ we certify is an idealization that degrades as $\Rwd$
departs from additivity. We have not characterized this degradation quantitatively; the gap between
idealized and approximate separability is a large theory-to-practice risk, and the empirical
separability of the embedding factors \citep{xiao2020whatnot,wang2020alignuniform,leace2023} is an
assumption our probes support only locally for content, and not at all for speaker.

\paragraph{OSA-3 is static geometry, not convergence; finite-time convergence is open.}
OSA-3 establishes two \emph{static} facts --- the rollout deficit $\beta\,\KL{\qz}{\qs}\ge0$
(\texttt{rollout\_deficit}) and the factorization of the optimum $\qs=\prod_i q^*_i$
(\texttt{qstar\_product}) --- and nothing dynamical. The operational update we actually run is JitRL's
\citep{jitrl2026} no-gradient credit assignment, whose guarantee is \emph{asymptotic under slow drift}
of the rollout distribution. The two regimes do not compose into a finite-time statement: we do
\emph{not} have a monotone-improvement guarantee for the gated update over the full, continuous,
only-approximately-separable space at a finite horizon. We read the $\beta$-KL term as a trust region
by analogy with proximal/trust-region dynamics, but this analogy is an open conjecture, not a theorem,
and it bears on the stability claims that agentic deployments care about most
\citep{ace2025,expfollowing2025,remembering2026}. We accordingly demote OSA-3 to a motivating static
geometry and treat convergence of the gated update as an open problem.

\paragraph{The isolated Hoeffding lemma and the T2 order-statistics \texttt{sorry} are load-bearing assumptions.}
Our machine-checked development is sorry-free in its qualitative core --- including the gain identity,
its nonnegativity, the flat-reward collapse, the partition-function and Gibbs-optimum product
factorizations, the strict per-non-degenerate-block positivity \texttt{gain\_pos\_of\_nonconstant},
and \texttt{rollout\_deficit} --- but two \emph{quantitative} results carry named, undischarged
hypotheses. The isolated Hoeffding lemma underwrites the $\spread^2/(8\beta)$ ceiling but its
concentration constant is assumed rather than derived from the actual reward's bounded-difference
structure; if the per-context reward is heavier-tailed than the assumed sub-Gaussian envelope, the
constants are optimistic. The T2 order-statistics result behind the best-of-$N$ spread $\spread$
retains a \texttt{sorry}: the asymptotics of the maximum order statistic
\citep{yang2024asymptotics,beirami2024bon} are invoked as a hypothesis, not proved in our
formalization. We flag both explicitly so that no reader mistakes ``machine-checked'' for
``unconditionally proved''; the checked artifact is sound \emph{relative to} these two assumptions.

\paragraph{Goodhart on verifiable rewards, and the sharper risk of \emph{model-dependent} rewards.}
Inference-time optimization against a fixed reward is, by construction, an over-optimization engine:
pushing $N$ or lowering $\beta$ moves probability mass toward whatever the reward measures, including
its defects \citep{gao2023overopt,skalse2022rewardhacking}. Verifiable speech rewards (WER, exact
match, intent accuracy) are more robust than learned reward models, but they are not immune, and
``spurious rewards'' can induce gains that do not reflect the intended capability
\citep{spuriousrewards2025,rlvrincentive2025}. A sharper, easily overlooked variant arises whenever
the reward is \emph{computed from the same frozen model} --- e.g.\ a probe read off the model's own
embedding, or a surrogate verifier instantiated from the model itself. Such a model-dependent reward
can be high-spread for reasons internal to the model rather than because it tracks the target
capability, so the certified $\gain$ may measure the model's self-consistency, not the skill. Because
we change no weights the failure is recoverable (lower $N$, raise $\beta$, decouple the verifier), but
the gate must be tuned against held-out, contamination-controlled splits
\citep{contaminationsurvey2024,gorman2019splits} with an \emph{externally} grounded reward, or the
gain is an artifact of the read-out.

\paragraph{Paralinguistic keys are dual-use: a mis-estimated label both mis-routes and mis-rewards.}
The same emotion/speaker estimate that \emph{routes} a query (selecting the context and skill) also
\emph{rewards} the rollout that produced it. A mis-estimated SER label therefore enters two coupled
loops with the same sign of error: it sends the query to the wrong context \emph{and} reinforces the
embedding that yielded the wrong label, a positive-feedback pathway that idealized separability hides.
This is sharpened by evidence that omni models tend to ``read, not listen'' on paralinguistic cues
\citep{voxparadox2026,listenortranscribe2025} and that speaker identity resists in-context correction
\citep{spkverifllm2026} --- exactly the regime where our speaker probe sits at chance. The mitigation
--- decoupling the routing estimator from the reward estimator --- is future work, not a property of
the current gate.

\paragraph{The two-omni ``vector-omni memory'' collapses for paralinguistics; we reposition the design.}
A natural reading of the system is two frozen omni models --- one as policy (Operator-B), one as a
vector memory whose embeddings index past episodes (Operator-A) --- with the omni memory carrying
paralinguistic keys. Our speaker result refutes the load-bearing version of that design: at
$\approx 0.04$ accuracy (near chance) across all layers, the omni embedding does \emph{not} support
paralinguistic retrieval, so a memory keyed on it would route on noise. We therefore reposition the
contribution honestly: the omni embedding is reserved for \emph{content} retrieval, where its spread
is real, and paralinguistic memory is delegated to classic, purpose-built speech encoders
(ECAPA-style speaker embeddings, a dedicated SER head) rather than to the omni read-out. The
``two-omni'' framing is downgraded to ``omni policy $+$ omni content memory $+$ classical
paralinguistic encoders.''

\paragraph{Memory poisoning via acoustic spoofing.}
When the optimized embedding is written to agent memory, the acoustic channel becomes an injection
surface: a spoofed speaker or emotion can plant a poisoned key that later retrieval amplifies
\citep{agentpoison2024,mempoison2026}. Unlike text poisoning, the attack needs no lexical trigger ---
it rides paralinguistic features the model already over-trusts. We provide no audit defense here;
memory hygiene and provenance \citep{anatomymem2026} are required before the embedding-as-memory-key
design is safe to deploy.

\paragraph{Small-sample intervals and provenance discrepancies under reconciliation.}
Several preliminary results are reported on modest evaluation sets, where bootstrap intervals
\citep{bisani2004bootstrap} are wide and statistical power is low
\citep{card2020power,dror2018hitchhiker}; effects near the resolution of a few hundred utterances
should be read as directional, not confirmed, and adequately powered, seed-averaged
\citep{henderson2018drlmatters} replication on the frozen, contamination-controlled split is a
precondition for any quantitative headline. Relatedly, the emotion re-run above already surfaced a
provenance discrepancy between an earlier single-seed, oracle-test-layer number and the multi-seed,
dev-selected protocol; we are reconciling all reported figures to the committed artifacts and the
dev-selected, paired-delta protocol, and any residual mismatch between an in-text number and its
artifact should be read as the artifact governing.

\section{Conclusion}
\label{sec:conclusion}

We have argued, and for its qualitative core machine-checked, a single reframed thesis: the
optimization gain available to training-free RL --- reward-guided, inference-time search with no
gradient and no weight or structure change --- is governed by the \emph{realized reward spread} of the
(action-space, reward) pair, a property of the \emph{reward structure} and not of model class. This is
the proven core, OSA-1: the gain is the divergence $\gain=\beta\,\KL{\qo}{\qs}$ between what the model
already does and what the reward would have it do, sandwiched as $0\le\gain\le\spread^2/(8\beta)$. A
flat reward --- e.g.\ instruction-conditioning of a frozen embedding read-out --- yields vanishing
gain; a high-spread reward --- e.g.\ task-level verifiable selection --- yields large gain \emph{even
on a single frozen model}, as our own single-model Operator-B results show (SLURP $+0.330$, URO
$+0.335$). The lever is the spread, not the agent.

OSA-2 sharpens this for composed rewards. We prove that gain is \emph{additive} across
context-isolated blocks (\texttt{gain\_product}) and that each \emph{non-degenerate} block contributes
\emph{strictly} positive gain (\texttt{gain\_pos\_of\_nonconstant}, sorry-free: a non-constant reward
forces $\gain>0$). But we also prove (\texttt{qstar\_product}) that the isolated optimum equals the
monolithic optimum on the product reward $\Rwd=\sum_i\Rwd_i$ --- so context isolation \emph{per se}
adds nothing; the only new gain comes from introducing additional non-degenerate verifiable rewards.
Whether a \emph{frozen speech model} actually affords non-degenerate per-block rewards is an empirical
question (\Cref{conj:spreadfloor}), not a theorem.

On that empirical question our own data return a central, honest negative for the flagship's
paralinguistic factors. A rigorous multi-seed re-run of the emotion-pooling gain (dev-selected layer,
paired-delta bootstrap, seeds $42,7,123,2024,31337$) gives mean $\Delta=+0.037$ with the paired CI
excluding zero in only $2$ of $5$ seeds --- the earlier $+0.097$ was the best seed under oracle
test-layer selection, a selection leak --- and the speaker probe is at chance ($\approx 0.04$) across
all $37$ layers. So the agentic-recovery move is empirically supported for \emph{content/intent} but,
on our evidence, likely \emph{fails} for paralinguistics. We report this not as a setback but as a
result: it scopes the thesis to where it has teeth and refutes the load-bearing version of a
two-omni paralinguistic memory, which we accordingly reposition as an omni policy plus omni
\emph{content} memory plus classical speaker/SER encoders.

OSA-3 we demote to motivating static geometry: it exhibits a positive static rollout deficit and a
static factorization of the optimum, but a finite-time monotone-convergence guarantee for the gated,
credit-assigned update over the full, only-approximately-separable space remains \emph{open}. The
trust-region reading of the $\beta$-KL term is an analogy we conjecture, not a theorem we prove.

We close by reaffirming the discipline that makes the claim falsifiable: a survey-first posture that
treats our own moat as an absence of evidence to be defended, not assumed; a commitment to
\emph{externally verifiable} rewards over learned or model-internal proxies, precisely to blunt the
Goodhart pressure that inference-time optimization invites; and a no-gradient, no-weight-change
constraint that keeps every failure recoverable at inference rather than frozen into parameters. The
promise of \emph{activating}, rather than retraining, the latent competence of frozen speech models is
real and near-term for content and intent; for paralinguistics it is, on present evidence, an open and
likely negative hypothesis --- and we state the whole only to the precision that the proofs and the
measurements can currently bear.